\begin{solucion}
    Dado que el multigrado de $f$ es $(4,0,0)$,  entonces construimos el polinomio:
    \begin{alignat*}{2}
      g_1 &=&&e_1^{4-0}e_2^{0-0}e_3^{0-0}\\
        &=&&e^{4}\\
        &=\bigl(x&&+y+z\bigr)^{4}\\
        &=(x&&+y)^4+4(x+y)^3z + 6(x+y)^2z^2 + 4(x+y)z^3 + z^4\\
        &=x^{4} &&+4x^{3}y + 6x^{2}y^{2} + 4xy^{3}+ y^4 + 4\bigl(x^{3}+3x^{2}y+3xy^{2}+y^{3}\bigr)z + 6\bigl(x^{2}+2xy+y^{2}\bigr)z^{2} \\
        &\phantom{=}&&+ 4(x+y)z^{3} + z^{4}\\
    &=\text{\setulcolor{red}\ul{$x^{4}$}}
      &&+\text{\setulcolor{blue}\ul{$4x^{3}y$}}
       +\text{\setulcolor{green}\ul{$6x^{2}y^{2}$}}
       +\text{\setulcolor{blue}\ul{$4xy^{3}$}}
       +\text{\setulcolor{red}\ul{$y^{4}$}}+\text{\setulcolor{blue}\ul{$4x^{3}z$}}
       +\text{\setulcolor{orange}\ul{$12x^{2}yz$}}
       +\text{\setulcolor{orange}\ul{$12xy^{2}z$}}
       +\text{\setulcolor{blue}\ul{$4y^{3}z$}}
       +\text{\setulcolor{green}\ul{$6y^{2}z^{2}$}} \\
    &\phantom{=}&&+\text{\setulcolor{green}\ul{$6x^{2}z^{2}$}}
       +\text{\setulcolor{orange}\ul{$12xyz^{2}$}}
       +\text{\setulcolor{blue}\ul{$4xz^{3}$}}
       +\text{\setulcolor{blue}\ul{$4yz^{3}$}}
       +\text{\setulcolor{red}\ul{$z^{4}$}} \\
    &=&&\text{\setulcolor{red}\ul{$x^{4} + y^{4} + z^{4}$}}
      +\text{\setulcolor{blue}\ul{$4(x^{3}y + x^{3}z + xy^{3} + y^{3}z + xz^{3} + yz^{3})$}}
       +\text{\setulcolor{green}\ul{$6(x^{2}y^{2} + x^{2}z^{2} + y^{2}z^{2})$}} \\
    &\phantom{=}&&+\text{\setulcolor{orange}\ul{$12(x^{2}yz + xy^{2}z + xyz^{2})$}} \\
    \end{alignat*}
    Por lo tanto:
    \begin{alignat*}{2}
        f_{1}&=f&&-g_1\\
        &=(x^{4}&&+y^{4}+z^{4})-x^{4} + y^{4} + z^{4}
        +4(x^{3}y + x^{3}z + xy^{3} + y^{3}z + xz^{3} + yz^{3})\\
        &\phantom{=}&&+6(x^{2}y^{2} + x^{2}z^{2} + y^{2}z^{2}) +12(x^{2}yz + xy^{2}z + xyz^{2}) \\
        &=&&-4(x^{3}y + x^{3}z + xy^{3} + y^{3}z + xz^{3} + yz^{3})-6(x^{2}y^{2} + x^{2}z^{2} + y^{2}z^{2})\\
        & &&-12(x^{2}yz + xy^{2}z + xyz^{2}) \\
    \end{alignat*}

    Por lo que
De esta manera, tenemos que $\mdeg(f_{1})=(3,1,0)$ por lo construimos:
\begin{alignat*}{2}
    g_2 &=&&e_1^{3-1}e_2^{1-0}e_3^{0-0}\\
        &=e_1^{2}&&e_2^1\\
        &=\bigl(x&&+y+z\bigr)^2\bigl(xy +yz+zx\bigr)\\
        &=\bigl((x&&+y)^2 + 2z(x+y) + z^2\bigr)\bigl(xy +yz+zx\bigr)\\
        &=\bigl(x^2 &&+ 2xy + y^2 + 2zx + 2zy + z^2\bigr)\bigl(xy +yz+zx\bigr)\\
        &=x^{3}y
          &&+2x^{2}y^{2}
          +xy^{3}
          +\text{\setulcolor{magenta}\ul{$2x^{2}yz$}}
          +\text{\setulcolor{cyan}\ul{$2xy^{2}z$}}
          +\text{\setulcolor{orange}\ul{$xyz^{2}$}}\\
            & &&
          +\text{\setulcolor{magenta}\ul{$x^{2}yz$}}
          +\text{\setulcolor{cyan}\ul{$2xy^{2}z$}}
          +y^{3}z
          +\text{\setulcolor{orange}\ul{$2xyz^{2}$}}
          +2y^{2}z^{2}
          +yz^{3}\\
            & &&
          +\text{\setulcolor{magenta}\ul{$2x^{2}yz$}}
          +\text{\setulcolor{cyan}\ul{$xy^{2}z$}}
          +x^{3}z
          +2x^{2}z^{2}
          +\text{\setulcolor{orange}\ul{$2xyz^{2}$}}
          +xz^{3}\\
        &=
          \text{\setulcolor{red}\ul{$x^{3}y$}}
          &&+\text{\setulcolor{blue}\ul{$2x^{2}y^{2}$}}
          +\text{\setulcolor{red}\ul{$xy^{3}$}}
          +\text{\setulcolor{green}\ul{$5x^{2}yz$}}
          +\text{\setulcolor{green}\ul{$5xy^{2}z$}}
          +\text{\setulcolor{green}\ul{$5xyz^{2}$}}
          +\text{\setulcolor{red}\ul{$y^{3}z$}}
          +\text{\setulcolor{blue}\ul{$2y^{2}z^{2}$}}\\
          & &&+\text{\setulcolor{red}\ul{$yz^{3}$}}
          +\text{\setulcolor{red}\ul{$x^3z$}}
          +\text{\setulcolor{blue}\ul{$2x^2z^2$}}
          +\text{\setulcolor{red}\ul{$xz^{3}$}}\\
        &=&&\text{\setulcolor{red}\ul{$x^{3}y+ xy^{3} + y^{3}z + x^{3}z + xz^{3} + yz^{3}$}}
        +\text{\setulcolor{blue}\ul{$2\bigl(x^{2}y^{2} + y^{2}z^{2}+x^{2}z^{2}\bigr)$}}\\
        & &&+\text{\setulcolor{green}\ul{$5\bigl(x^{2}yz + xy^{2}z + xyz^{2}\bigr)$}}
    \end{alignat*}

    Por lo tanto:

    \begin{alignat*}{2}
        f_{2}&=&&f_{1}-(-4g_{2})\\
        &=&&-\cancel{4(x^{3}y+ xy^{3} + y^{3}z + x^{3}z + xz^{3} + yz^{3})} - 6(x^{2}y^{2} + x^{2}z^{2} + y^{2}z^{2})\\
        & &&-12(x^{2}yz + xy^{2}z + xyz^{2})+\cancel{4\bigl(x^{3}y + xy^{3} + y^{3}z + x^{3}z + xz^{3} + yz^{3}\bigr)}\\
        & && + 8\bigl(x^{2}y^{2} + y^{2}z^{2}+x^{2}z^{2}\bigr)+20\bigl(x^{2}yz + xy^{2}z + xyz^{2}\bigr)\\
        &=&&2(x^{2}y^{2} + x^{2}z^{2} + y^{2}z^{2}) + 8(x^{2}yz + xy^{2}z + xyz^{2})\\
    \end{alignat*}
    Vemos que $\mdeg(f_{2})=(2,2,0)$ por lo que construimos:
    \begin{alignat*}{2}
        g_3 &=e_1^{2-2}e_2^{2-0}e_3^{0-0}\\
            &=e_2^{2}\\
            &=\bigl(xy + yz + zx\bigr)^{2}\\
            &=(xy + yz)^2 + 2(xy + yz)xz + x^2z^2\\
            &=\text{\setulcolor{red}\ul{$x^2y^2$}}
              +\text{\setulcolor{blue}\ul{$2xy^2z$}}
              +\text{\setulcolor{green}\ul{$y^2z^2$}}
              +\text{\setulcolor{blue}\ul{$2x^2yz$}}
              +\text{\setulcolor{green}\ul{$2xyz^2$}}
              +\text{\setulcolor{red}\ul{$x^2z^2$}}\\
            &= \text{\setulcolor{red}\ul{$x^2y^2 + y^2z^2 + x^2z^2$}}
              +\text{\setulcolor{blue}\ul{$2(x^2yz + xy^2z + xyz^2)$}}\\
    \end{alignat*}
    Luego:
    \begin{alignat*}{2}
        f_{3}&=&&f_{2}-2g_{3}\\
        &=&&\cancel{2(x^{2}y^{2} + x^{2}z^{2} + y^{2}z^{2})} + 8(x^{2}yz + xy^{2}z + xyz^{2})-\cancel{2(x^{2}y^{2} + y^{2}z^{2} + x^{2}z^{2})}\\
        & && - 4\bigl(x^{2}yz + xy^{2}z + xyz^{2}\bigr)\\
        &=&&4(x^{2}yz + xy^{2}z + xyz^{2})\\
    \end{alignat*}
    De esta manera, tenemos que $\mdeg(f_{3})=(2,1,1)$ por lo que construimos:
    \begin{alignat*}{2}
        g_4 &=e_1^{2-1}e_2^{1-1}e_3^{1-0}\\
            &=e_1e_3\\
            &=\bigl(x+y+z\bigr)(xyz)\\
            &=x^{2}yz + xy^{2}z + xyz^{2}\\
    \end{alignat*}
    Por lo que:
    \begin{alignat*}{2}
        f_{4}&=f_{3}-4g_{4}\\
        &=4\bigl(x^{2}yz + xy^{2}z + xyz^{2}\bigr)-4\bigl(x^{2}yz + xy^{2}z + xyz^{2}\bigr)\\
        &=0\\
    \end{alignat*}
    Por lo tanto, el polinomio $f$ es reducible y se puede escribir como:
    \begin{alignat*}{2}
        f &=g_1 -4g_2 + 2g_3 + 4g_4\\
        &=e_1^{4} - 4e_1^2e_2 + 2e_2^2 + 4e_1e_3\\
    \end{alignat*}
    $$
    f
      = e_{1}^{\,4}\,-\,4e_{1}^{\,2}e_{2}\,+\,2e_{2}^{\,2}\,+\,4e_{1}e_{3}.$$
\end{solucion}
